\documentclass[UTF8,fontset=none,11pt,a4paper]{ctexart}
\usepackage{ipara-handbook}
\usepackage{graphicx}
\handbooksetup{DS-12 外部排序}{408 · 数据结构 DS-12}{Block I/O · Runs · Multiway Merge · Passes}
\begin{document}
\handbooktitle{外部排序}{数据不能全进内存时，怎样减少块 I/O 轮数}{408 · 数据结构 Topic DS-12}

\begin{corebox}{母问题：当 I/O 远贵于比较，排序模型怎样改写？}
外部排序的生成链是
\[
\text{Block-I/O Cost}\to\text{Initial Runs}\to\text{Multiway Merge}\to\text{Pass Count}\to\text{I/O Tradeoff}.
\]
数据先分成能装入内存的初始归并段并在内存中排序，再用 $k$ 路归并逐轮合并。总成本主要由读写块次数和归并轮数决定，不能直接套内部排序的比较次数。
\end{corebox}
\begin{methodbox}{本册定位与 Owner 边界}
\begin{itemize}
\item \kw{Owns：}初始归并段、置换选择、多路归并、败者树、缓冲区、I/O 轮数与终止条件。
\item \kw{Uses：}Atlas Foundation 的 cost vector；存储系统的 block 概念只作为外部接口，不复制文件系统机制。
\item \kw{Stops：}内存内局部排序归 DS11；设备调度、缓存一致性和文件系统完整语义不在本册。
\end{itemize}
\end{methodbox}
\tableofcontents\newpage

\section{初始段与多路归并}
设内存可容纳 $M$ 条记录，先读入至多 $M$ 条并排序，写成一个 run。归并时每个输入 run 保留一个当前首元素，取最小者输出并补充同一路的下一条；候选结构可用最小堆或败者树。若有 $r$ 个初始段、每轮 $k$ 路，轮数约为 $\lceil\log_k r\rceil$，每轮都要读写全部数据。
普通方法每次只处理能装入内存的 $M$ 条，所以初始段数约为 $r=\lceil N/M\rceil$。置换--选择排序把“选择输出”和“从输入补位”交错进行：下一条必须选择所有 key 不小于上一条输出值中的最小者；小于该值的记录冻结到下一 run。随机输入下平均 run 长度约为 $2M$，因而初始段数大致减半，但生成阶段需要一个可维护候选极值的结构，成本约为 $O(N\log M)$。

置换--选择实现中，内存槽位至少有“活动”和“冻结”两种状态。输出当前最小活动记录后读入新记录：若新 key 不小于刚输出值，继续留在本 run；否则标记冻结，等待 run 结束。活动槽位耗尽时关闭当前 run、把冻结集合解冻并开始下一 run。若没有把“冻结”与“真正输入耗尽”分开，容易提前结束或把下降 key 错写进当前有序段；若输入已经读完，剩余活动槽仍需全部输出。
\begin{examplebox}{为什么败者树值得出现}
多路归并每输出一条记录都要更新“哪一路最小”。普通堆和败者树都维护候选极值；败者树显式保存比较路径，便于重复更新同一路。二者都不改变多路归并的正确性不变量。
\end{examplebox}

\subsection{多路归并、缓冲区与败者树}
朴素地从 $k$ 个当前首元素中找最小值，每输出一条记录要做 $k-1$ 次比较；$k$ 增大虽然减少归并趟数，却可能让内部比较抵消 I/O 收益。败者树在内部节点保存每场比较的败者，让胜者沿路径晋级；冠军输出后只需沿该路回到根，单次调整约为 $\lceil\log_2 k\rceil$ 次比较，内部比较成本与 $k$ 的线性增长脱钩。相比胜者树，败者树保存的正是新元素沿路径再次会遇到的对手，更新时不必重新访问兄弟节点。

\begin{center}
\providecolor{themebg}{HTML}{FAFAF7}
\providecolor{themefg}{HTML}{111111}
\providecolor{themegray}{HTML}{666666}
\providecolor{themecurve}{HTML}{1D63B8}
\providecolor{themealert}{HTML}{C53030}
\providecolor{themeamber}{HTML}{B86E00}
\begin{tikzpicture}[>=Stealth,font=\small,color=themefg,text=themefg]
\tikzset{
  leaf/.style={rounded corners=3pt,draw=themefg,fill=themebg,minimum width=1.8cm,minimum height=.72cm,align=center},
  loser/.style={circle,draw=themeamber,fill=themeamber!10!themebg,minimum size=8.5mm,inner sep=0pt},
  winner/.style={rounded corners=3pt,draw=themecurve,fill=themecurve!10!themebg,minimum width=2.35cm,minimum height=.78cm,align=center},
  edge/.style={themegray,line width=.9pt},
  path/.style={themecurve,line width=1.6pt},
  note/.style={rounded corners=4pt,draw=themegray,fill=themefg!3!themebg,inner xsep=7pt,inner ysep=6pt,align=left}
}
\node[anchor=west,font=\bfseries\large] at (0,7.0) {4 路败者树：每次只让“刚补入的那一路”沿叶到根重新比赛};
\node[anchor=west,text=themegray] at (0,6.45)
  {叶子只保存各 run 的当前首元素；内部节点保存本场败者，顶部另保留全局胜者。};

% initial tree
\node[font=\bfseries,text=themecurve] at (4.0,5.55) {输出前：当前首元素 $3,6,2,8$};
\node[winner] (w0) at (4.0,4.85) {全局胜者 $R_2:2$};
\node[loser] (lroot) at (4.0,3.75) {3};
\node[loser] (ll) at (2.45,2.65) {6};
\node[loser] (lr) at (5.55,2.65) {8};
\node[leaf] (r0) at (1.35,1.35) {$R_0$\\3};
\node[leaf] (r1) at (3.55,1.35) {$R_1$\\6};
\node[leaf,draw=themecurve] (r2) at (5.75,1.35) {$R_2$\\2};
\node[leaf] (r3) at (7.95,1.35) {$R_3$\\8};
\draw[edge] (r0)--(ll); \draw[edge] (r1)--(ll); \draw[edge] (r2)--(lr); \draw[edge] (r3)--(lr);
\draw[edge] (ll)--(lroot); \draw[edge] (lr)--(lroot); \draw[edge] (lroot)--(w0);
\node[anchor=west,font=\scriptsize,text=themegray,align=left] at (8.45,3.2)
  {败者：$6,8,3$\\胜者：$2$};

\draw[themegray,line width=.45pt] (9.65,.65)--(9.65,5.75);

% update tree
\node[font=\bfseries,text=themecurve] at (13.0,5.55) {输出 2 后，$R_2$ 补入 7};
\node[winner] (w1) at (13.0,4.85) {新胜者 $R_0:3$};
\node[loser,draw=themecurve] (uroot) at (13.0,3.75) {7};
\node[loser] (ul) at (11.45,2.65) {6};
\node[loser,draw=themecurve] (ur) at (14.55,2.65) {8};
\node[leaf] (u0) at (10.35,1.35) {$R_0$\\3};
\node[leaf] (u1) at (12.25,1.35) {$R_1$\\6};
\node[leaf,draw=themecurve] (u2) at (14.15,1.35) {$R_2$\\7};
\node[leaf] (u3) at (16.05,1.35) {$R_3$\\8};
\draw[edge] (u0)--(ul); \draw[edge] (u1)--(ul); \draw[path] (u2)--(ur); \draw[edge] (u3)--(ur);
\draw[edge] (ul)--(uroot); \draw[path] (ur)--(uroot); \draw[path] (uroot)--(w1);
\node[anchor=west,font=\scriptsize,text=themecurve] at (9.85,.5)
  {蓝线：$7$ 只沿原比较路径重赛；左侧 $R_0/R_1$ 子树不动。};

\node[note,anchor=west,text width=7.9cm] at (0,-.2)
  {\textbf{复杂度来源：}一次输出只改变一个 run 的当前首元素，因此只有该叶到根的比较链失效。树高为 $\lceil\log_2 k\rceil$，单次选择更新为 $O(\log k)$，但归并 I/O 趟数和缓冲区需求不因此改变。};
\end{tikzpicture}

\end{center}
图中内部节点保存当前比较的败者，顶部单独保存全局胜者；输出某一路并补入新值后，只有该叶到根的比较链失效，其他子树不需要重赛。

每一路至少需要输入缓冲区，输出还需一个输出缓冲区；$k$ 不能脱离内存预算盲目增大。某一路耗尽时可把它的当前值替换为 $+\infty$，当冠军也是 $+\infty$ 才表示全部输入耗尽。初始化败者树可先让内部节点指向 $-\infty$ 虚拟叶，再逐路插入首元素，避免未初始化比较。

缓冲区设计还要区分“记录耗尽”和“缓冲块耗尽”：前者把该路从候选集合移除，后者只触发一次块读入。若采用 $k$ 路归并而内存只有 $B$ 个缓冲块，通常需为每路输入和输出至少保留一个块，故 $k$ 受 $B$ 约束；增加路数虽可降低理论趟数，却可能使每路缓冲过小、I/O 次数和管理开销上升。题目给出具体缓冲预算时，应先算可行的 $k$ 再算 $\lceil\log_k r\rceil$。

\subsection{归并趟数、最佳归并树与虚段}
若初始段数为 $r$、采用平衡 $k$ 路归并，段数变化为 $r,\lceil r/k\rceil,\lceil r/k^2\rceil,\ldots$，所以趟数为 $S=\lceil\log_k r\rceil$；每趟读一遍写一遍，I/O 量约为 $2NS$。初始段长度不等时，先合并哪些段会影响它们被重复读写的次数。把段长作为叶权，归并树的总读写量等于 $2\times WPL$；每次取最小的 $k$ 个权值合并，得到 $k$ 叉哈夫曼式最佳归并树。

严格 $k$ 叉树要求
\[
\frac{r-1}{k-1}\in\mathbb Z.
\]
令 $u=(r-1)\bmod(k-1)$。若 $u\ne0$，需补 $k-u-1$ 个权值为零的虚段，等价于第一次只做 $u+1$ 路归并；虚段应位于树的最底层，不能拖到最后，否则会让真实段多经历一趟 I/O。虚段只参与归并树计算，不是外存中的真实记录。

\section{I/O 不变量与终止}
每个 run 的当前指针只向前移动；输出序列始终是所有当前首元素中的最小值；某一路耗尽后从候选集合移除；所有输入耗尽才结束。缓冲区满/空是 I/O 事件，不是记录顺序事件，必须分开书写。块大小、内存缓冲数量和归并路数改变的是轮数与常数，不改变“每一路有序”的前提。
外部排序的总时间应拆成初始段内部排序、外存读写和归并阶段内部处理三部分。在 408 的块 I/O 模型中通常假设 $t_{IO}\gg t_{cmp}$，所以优化第一目标是减少归并趟数和总写回量；但扩大 $k$ 会消耗更多输入缓冲，可能降低每路缓冲效率。真实 workload 若包含 JSON 解析、复杂比较、压缩或校验，CPU 处理可能反而主导，必须测量后再定优化对象。题目给出块大小时，应按“块读/块写”计数，不把每条记录比较次数当作 I/O 次数。

\section{范围分桶：条件满足时绕开通用归并}
若 key 的全局范围已知，可把输入按互不重叠且有序的区间写入多个外存桶；只要每桶能在内存中排序，最后按桶区间顺序串接即可得到全局有序结果。它省去多轮通用归并，但正确性依赖三项前提：桶间范围严格有序；边界值只进入唯一桶；每桶规模不超过可处理预算。

数据倾斜会让单个热点桶溢出，此时需要二次分桶、采样后重设边界，或退回外部归并。按 Hash 值分区不能用于全局排序，因为不同 Hash 桶没有 key 顺序；它适合让相同 key 落到同一分区，从而处理去重、频次或集合交集，这些组合路线由 DS-I01 统筹。范围桶也不同于 DS11 的内存桶排序：这里的核心代价是分区写入、桶内处理与最终顺序读取的 I/O。

\section{瓶颈由 workload 证据决定}
统一成本模型可写为
\[
T=T_{read}+T_{parse}+T_{compare}+T_{structure}+T_{write}+T_{merge}.
\]
纯整数 Top-K 常常表现为读写主导；结构化记录外排可能因解析和比较让 CPU 主导。两种观察并不矛盾，它们说明“数据大”不足以决定优化方向。设计前应分别测吞吐、CPU 利用率、峰值内存、临时文件体积与各阶段耗时，再判断应增加缓冲、减少趟数、简化 key 提取，还是并行化某一阶段。

并发 worker 会同时放大输入缓冲、输出缓冲和解析对象的内存占用；串行归并与最终写回还会受 Amdahl 上限约束。线程数从 $p$ 增到 $2p$ 不保证吞吐翻倍，尤其在设备带宽饱和、垃圾回收或共享队列竞争时。以上是工程 Extension，不改变 408 题设下按块和趟数计数的 Canonical 模型，但防止把该模型误当成所有真实系统的测量结论。

\section{代码与测试证据}
完整实现位于 \codepath{code/ds12_external_sort.hpp}，测试位于 \codepath{code/ds12_external_sort_test.cpp}。可执行核心用内存中的 run 模拟外部边界，输出合并结果并计算所需轮数。
\lstinputlisting[language=C++,firstline=10,lastline=82,firstnumber=1,numbers=left,caption={有序 runs 的多路归并与轮数估算}]{30_408/10_数据结构/12_外部排序/code/ds12_external_sort.hpp}
\lstinputlisting[language=C++,firstline=8,lastline=40,firstnumber=1,numbers=left,caption={空 run、单 run 与多路耗尽测试}]{30_408/10_数据结构/12_外部排序/code/ds12_external_sort_test.cpp}
\begin{lstlisting}[language=bash]
clang++ -std=c++17 -Wall -Wextra -Werror -pedantic code/ds12_external_sort_test.cpp -o /tmp/ds12_test
/tmp/ds12_test
\end{lstlisting}

\section{复杂度与概念边界}
\begin{center}\small
\begin{tabularx}{\linewidth}{L{2.8cm}C{1.8cm}C{1.8cm}YY}\toprule
阶段 & 比较/更新 & I/O 视角 & 不变量 & 边界\\\midrule
初始段 & 每段内存排序 & 读写 $N$ & 每段有序 & $N\le M$\\
$k$ 路归并 & 每条更新候选 & 每轮读写 $N$ & 各路首元素候选 & 空 run\\
总过程 & 依赖实现 & $O(N\log_k r)$ 轮次量级 & 指针单调前进 & 输入耗尽\\\bottomrule\end{tabularx}\end{center}

\section{做题控制协议}
先标 $N,M,k,r$，再算初始段数量和归并轮数；每一轮分别写“读、比较、输出、补位、写”。看到全部数据能进内存，停止套外部 I/O 公式并转 DS11；看到块大小或缓冲区变化，重新计算轮数而非只比较关键字次数。
\begin{corebox}{压缩复原}
五问：内存一次容纳多少？初始 run 怎样形成？当前候选来自几路？哪一路耗尽如何退出？总成本按比较还是块 I/O 计？
\end{corebox}
\end{document}
